\chapter{Monero 交易相关知识证明 (Monero Transaction-Related Knowledge Proofs)}
\label{chapter:tx-knowledge-proofs}
%original draft by Sarang Noether, Ph.D.

\iffalse
https://github.com/monero-project/monero/pull/6329/files

https://monero.stackexchange.com/questions/8122/what-is-the-spendproofv1-or-outproofv1-in-the-details-of-a-sent-transa

https://monero.stackexchange.com/questions/9991/how-does-the-get-reserve-proof-command-work

https://github.com/monero-project/research-lab/issues/68
\fi

Monero 是一种货币，如同任何货币一样，其用途是复杂的。从企业财务、市场交易到法律仲裁，不同的利益相关方可能希望了解交易的详细信息。

你如何确定收到的钱来自特定的人？或者在对方否认的情况下，证明你确实向某人发送了某个 output 或交易？在 Monero 的公共账本中，发送者和接收者都是模糊的。你如何在不泄露私钥的情况下证明你拥有一定数量的资金？在 Monero 中，金额对观察者来说是完全隐藏的。

我们考虑了几种类型的交易断言，其中一些已在 Monero 中实现，可通过内置钱包工具使用。我们还概述了一个审计个人或组织所持全部余额的框架，该框架不会泄露其未来交易的信息。



\section{Monero 中的交易证明 (Transaction proofs in Monero)}
\label{sec:proofs-monero-proofs}

Monero 的交易证明正在进行更新 \cite{sarang-txproofs-updates-issue}。目前已实现的证明都是「版本 1」，不包含域分离。我们仅描述最先进的证明，无论它们是目前已实现的、计划在未来版本中实现的 \cite{sarang-txproofs-v2-update-pr}，还是可能会也可能不会被实现的假设性证明（第 \ref{subsec:proofs-owned-output-spent-unspentproof} 节 \cite{unspent-proof-issue-68} 和第 \ref{subsec:proofs-address-subaddress-correspond-subaddressproof} 节）。


\subsection{多基 Monero 交易证明 (Multi-base Monero transaction proofs)}
\label{subsec:proofs-multi-base-monero}

有一些细节需要注意。大多数 Monero 交易证明涉及多基证明（回忆第 \ref{sec:proofs-discrete-logarithm-multiple-bases} 节）。在相关之处，域分隔符为 $T_{txprf2} = \mathcal{H}_n(``\textrm{TXPROOF\_V2}")$。\footnote{正如第 \ref{sec:CLSAG} 节所述，哈希函数应通过前缀标签进行域分离。当前 Monero 交易证明的实现没有域分离，因此本章中的所有标签都是{\em 尚未}实现的功能。}被签名的消息\marginnote{src/wallet/ wallet2.cpp {\tt get\_tx\_ proof()}}通常（除非另有说明）为 $\mathfrak{m} = \mathcal{H}_n(\texttt{tx\_hash, \texttt{message}})$，其中 {\tt tx\_hash} 是相关交易的 ID（第 \ref{subsec:transaction-id} 节），{\tt message} 是证明者或第三方可以提供的可选消息，以确保证明者确实制作了证明而非窃取了他人的。

证明以 base-58 编码，这是一种最初为 Bitcoin 引入的二进制到文本的编码方案 \cite{base-58-encoding}。验证这些证明始终涉及先将其从 base-58 解码回二进制。注意验证者还需要访问区块链，以便使用交易 ID 引用来获取一次性地址等信息。\footnote{交易 ID 通常与证明分开传达。}\\

证明中密钥前缀的结构有些不对称，部分原因是累积的更新尚未重新组织。2 基「版本 2」证明的挑战按以下格式组装，如果「基密钥 1」是 $G$，则其在挑战中的位置用 32 个零字节填充，\vspace{.175cm}
\[c = \mathcal{H}_n(\mathfrak{m}\textrm{, public key 2, proof part 1, proof part 2, $T_{txprf2}$, public key 1, base key 2, base key 1})\]


\subsection{证明交易 input 的创建 ({\tt SpendProofV1})}
\label{subsec:proofs-input-creation-spendproof}

假设我们进行了一笔交易，并想要证明它。显然，通过在新消息上重新制作交易 input 的签名，任何验证者都只能得出结论：原始交易确实是我们制作的。重新制作交易{\em 所有} input 的签名意味着我们必须制作了整笔交易（回忆第 \ref{full-signature} 节），或者至少完全为其提供了资金。\footnote{正如我们将在第 \ref{chapter:txtangle} 章中看到的，制作了一个 input 签名的人不一定制作了所有 input 签名。}

所谓的「SpendProof」\marginnote{src/wallet/ wallet2.cpp {\tt get\_spend\_ proof()}}包含对交易所有 input 的重新制作的签名。重要的是，SpendProof 环签名重新使用了原始环成员，以避免通过环交集识别真正的签名者。

SpendProof 已在 Monero 中实现。为了编码传输给验证者，证明者将前缀字符串「{\tt SpendProofV1}」与签名列表连接起来。注意前缀字符串不在 base-58 编码中，也不需要编码/解码，因为其目的是为了人类可读性。

\subsubsection*{SpendProof}

SpendProof 出人意料地不使用 MLSAG，而是使用 Monero 最初的环签名方案 \marginnote{src/crypto/ crypto.cpp {\tt generate\_ ring\_signa- ture()}}，该方案用于最早的交易协议（RingCT 之前）\cite{cryptoNoteWhitePaper}。

\begin{enumerate}
	\item 计算 key image \(\tilde{K} = k^o_\pi \mathcal{H}_p(K^o_\pi)\)。

	\item 生成随机数 \(\alpha \in_R \mathbb{Z}_l\)，以及对 \(i \in \{1, 2, ..., n\}\) 但排除 \(i = \pi\) 的随机数 \(c_i, r_i \in_R \mathbb{Z}_l\)。

	\item 计算
	\[c_{tot} = \mathcal{H}_n(\mathfrak{m},[r_1 G + c_1 K^o_1],[r_1 \mathcal{H}_p(K^o_1) + c_1 \tilde{K}],...,[\alpha G],[\alpha \mathcal{H}_p(K^o_{\pi})],...,\textrm{etc.})\]

	\item 定义真正的挑战
	\[c_{\pi} = c_{tot} - \sum^{n}_{i=1,i\neq \pi} c_i\]

	\item 定义 \(r_{\pi} = \alpha - c_{\pi}*k^o_{\pi} \pmod l\)。
\end{enumerate}

签名为 $\sigma = (c_1, r_1,c_2,r_2,...,c_n,r_n)$。

\subsubsection*{验证 (Verification)}

验证\marginnote{src/wallet/ wallet2.cpp {\tt check\_spe- nd\_proof()}}给定交易上的 SpendProof 时，验证者使用相关引用交易中的信息（例如 key image，以及从其他交易获取一次性地址的 output 偏移量）确认所有环签名有效。

\begin{enumerate}
	\item 计算
	\[c_{tot} = \mathcal{H}_n(\mathfrak{m},[r_1 G + c_1 K^o_1],[r_1 \mathcal{H}_p(K^o_1) + c_1 \tilde{K}],...,[r_n G + c_n K^o_n],[r_n \mathcal{H}_p(K^o_n) + c_n \tilde{K}])\]

	\item 检查
	\[c_{tot} \stackrel{?}{=} \sum^{n}_{i=1} c_i\]
\end{enumerate}

\subsubsection*{原理 (Why it works)}

注意当只有一个环成员时，该方案与 bLSAG（第 \ref{blsag_note} 节）相同。添加伪造成员时，不是将挑战 $c_{\pi+1}$ 传入新的挑战哈希，而是将成员添加到原始哈希中。由于以下方程\vspace{.175cm}
\[c_{s} = c_{tot} - \sum^{n}_{i=1,i\neq s} c_i\]

对任意索引 $s$ 都平凡成立，验证者无法识别真正的挑战。此外，如果不知道 $k^o_{\pi}$，证明者将永远无法正确定义 $r_{\pi}$（概率可忽略不计的情况除外）。


\subsection{证明交易 output 的创建 ({\tt OutProofV2})}
\label{subsec:proofs-output-creator-outproof}

现在假设我们向某人发送了资金（一个 output），并想要证明它。交易 output 的核心包含三个组成部分：接收者地址、发送金额和交易私钥。金额是编码过的，因此我们实际上只需要地址和交易私钥即可开始。任何删除或丢失了交易私钥的人都将无法制作 OutProof，因此从这个意义上说，OutProof 是所有 Monero 交易证明中最不可靠的。\footnote{我们可以将「OutProof」理解为展示一个 output 从证明者处「发出」。相应的「InProof」（第 \ref{subsec:proofs-output-ownership-inproof} 节）展示的是「进入」证明者地址的 output。}

我们这里的任务是展示一次性地址是从接收者的地址生成的，并允许验证者重构 output 承诺。我们通过提供发送者-接收者共享秘密 $rK^v$ 来做到这一点，然后通过在基密钥 $G$ 和 $K^v$ 上签名一个 2 基签名（第 \ref{sec:proofs-discrete-logarithm-multiple-bases} 节）来证明我们创建了它，并且它与交易公钥和接收者地址相对应。验证者可以使用共享秘密来检查接收者\marginnote{src/wallet/ wallet2.cpp {\tt check\_tx\_ proof()}}（第 \ref{sec:one-time-addresses} 节）、解码金额（第 \ref{sec:pedersen_monero} 节），并重构 output 承诺（第 \ref{sec:pedersen_monero} 节）。我们提供普通地址和子地址的详细信息。

\subsubsection*{OutProof}

为向地址 $(K^{v},K^{s})$ 或子地址 $(K^{v,i},K^{s,i})$ 发送的 output 生成\marginnote{src/crypto/ crypto.cpp {\tt generate\_ tx\_proof()}}证明，交易私钥为 $r$，发送者-接收者共享秘密为 $rK^v$。回忆存储在交易数据中的交易公钥为 $rG$ 或 $rK^{s,i}$，取决于接收者是否使用子地址（第 \ref{sec:subaddresses} 节）。

\begin{enumerate}
	\item 生成随机数 $\alpha \in_R \mathbb{Z}_l$，并计算
	\begin{enumerate}
	    \item {\em 普通地址}：$\alpha G$ 和 $\alpha K^v$
	    \item {\em 子地址}：$\alpha K^{s,i}$ 和 $\alpha K^{v,i}$
	\end{enumerate}{}
	\item 计算挑战
	\begin{enumerate}
	    \item {\em 普通地址}：\footnote{此处的 `0' 值是 32 字节的零字节编码。}
	    \[c = \mathcal{H}_n(\mathfrak{m},[rK^v], [\alpha G], [\alpha K^v], [T_{txprf2}], [rG], [K^v], [0])\]
	    \item {\em 子地址}：
	    \[c = \mathcal{H}_n(\mathfrak{m},[rK^{v,i}], [\alpha K^{s,i}], [\alpha K^{v,i}], [T_{txprf2}], [rK^{s,i}], [K^{v,i}], [K^{s,i}])\]
	\end{enumerate}{}
	\item 定义响应\footnote{由于可用符号有限，我们遗憾地用 $r$ 同时表示响应和交易私钥。必要时将使用上标 `resp' 表示「response」以区分两者。} $r^{resp} = \alpha - c*r$。
	\item 签名为 $\sigma^{outproof} = (c, r^{resp})$。
\end{enumerate}{}

证明者\marginnote{src/wallet/ wallet2.cpp {\tt get\_tx\_ proof()}}可以生成一批 OutProof，然后一起发送给验证者。他将前缀字符串「{\tt OutProofV2}」与证明列表连接起来，其中每项（以 base-58 编码）包含发送者-接收者共享秘密 $r K^v$（子地址则为 $r K^{v,i}$）及其对应的 $\sigma^{outproof}$。我们假设验证者知道每条证明对应的地址。

\subsubsection*{验证 (Verification)}

\begin{enumerate}
    \item 计算挑战\marginnote{src/crypto/ crypto.cpp {\tt check\_tx\_ proof()}}
    \begin{enumerate}
        \item {\em 普通地址}：\vspace{.145cm}
	    \[c' = \mathcal{H}_n(\mathfrak{m},[rK^v], [r^{resp} G + c*r G], [r^{resp} K^v + c*r K^v], [T_{txprf2}], [rG], [K^v], [0])\]
	    \item {\em 子地址}：\vspace{.16cm}
	    \[c' = \mathcal{H}_n(\mathfrak{m},[rK^{v,i}], [r^{resp} K^{s,i} + c*r K^{s,i}], [r^{resp} K^{v,i} + c*r K^{v,i}], [T_{txprf2}], [rK^{s,i}], [K^{v,i}], [K^{s,i}])\]
    \end{enumerate}{}
    \item 若 $c = c'$，则证明者知道 $r$，且 $rK^v$ 是 $r G$ 和 $K^v$ 之间的合法共享秘密（概率可忽略不计的情况除外）。
    \item 验证者\marginnote{src/wallet/ wallet2.cpp {\tt check\_tx\_ key\_hel- per()}}应检查所提供的接收者地址是否能用于从相关交易生成一次性地址（普通地址和子地址的计算方式相同）
    \[K^s \stackrel{?}{=} K^o_t - \mathcal{H}_n(r K^v,t)\]
    \item 他们还应解码 output 金额 $b_t$，计算 output 掩码 $y_t$，并尝试重构对应的 output 承诺\footnote{有效的 OutProof 签名并不一定意味着所考虑的接收者就是真正的接收者。恶意证明者可以生成随机 view key $K'^v$，计算 $K'^s = K^o - \mathcal{H}_n(rK'^v,t)*G$，并提供 $(K'^v,K'^s)$ 作为名义接收者。通过重新计算 output 承诺，验证者可以更有信心地确认所讨论的接收者地址是合法的。然而，证明者和接收者可以合谋使用 $K'^v$ 编码 output 承诺，而一次性地址使用 $(K^v,K^s)$。由于接收者需要知道私钥 $k'^v$（假设 output 金额仍然可花费），这种程度的欺骗的实际用处令人质疑。为什么接收者不直接使用 $(K'^v,K'^s)$（或其他一次性地址）来接收整个 output 呢？由于 $C^b_t$ 的计算与接收者相关，我们认为所述的 OutProof 验证过程是足够的。换言之，证明者无法在不与接收者合谋的情况下用它来欺骗验证者。}\vspace{.175cm}
    \[C^b_t \stackrel{?}{=} y_t G + b_t H\]
\end{enumerate}{}


\subsection{证明 output 的所有权 ({\tt InProofV2})}
\label{subsec:proofs-output-ownership-inproof}

OutProof 展示证明者向某个地址发送了一个 output，而 InProof 展示某个地址收到了一个 output。它本质上是发送者-接收者共享秘密 $r K^v$ 的另一「面」。这次证明者证明他知道 $K^v$ 中的 $k^v$，并且结合交易公钥 $r G$ 后共享秘密 $k^v*r G$ 会出现。

一旦验证者有了 $r K^v$，他们就可以通过以下方式检查对应的一次性地址是否由证明者的地址所拥有\marginnote{src/wallet/ wallet2.cpp {\tt check\_tx\_ proof()}} $K^o - \mathcal{H}_n(k^v*rG,t)*G \stackrel{?}{=} K^s$（第 \ref{sec:multi_out_transactions} 节）。通过对区块链上所有交易公钥制作 InProof，证明者将暴露他所有的已拥有的 output。

直接给验证者 view key 会有同样的效果，但一旦他们拥有了该密钥，验证者将能够识别未来创建的 output 的所有权。使用 InProof，证明者可以保留对其私钥的控制权，代价是证明（然后验证）每个 output 是否被拥有所需的时间。

\subsubsection*{InProof}

InProof 的构建方式与 OutProof 相同\marginnote{src/crypto/ crypto.cpp {\tt generate\_ tx\_proof()}}, 只是基密钥现在是 $\mathcal{J} = \{G, r G\}$，公钥是 $\mathcal{K} = \{K^v, r K^{v}\}$，签名密钥是 $k^v$ 而非 $r$。我们将仅展示验证步骤以说明我们的意思。注意密钥前缀的顺序发生了变化（$r G$ 和 $K^v$ 互换了位置），以匹配每个密钥的不同角色。

多个 InProof 可以一起发送给验证者，这些 InProof 与同一地址拥有的多个 output 相关。\marginnote{src/wallet/ wallet2.cpp {\tt get\_tx\_ proof()}} 它们以字符串「{\tt InProofV2}」为前缀，每项（以 base-58 编码）包含发送者-接收者共享秘密 $r K^v$（或 $r K^{v,i}$）及其对应的 $\sigma^{inproof}$。

\subsubsection*{验证 (Verification)}

\begin{enumerate}
    \item 计算挑战\marginnote{src/crypto/ crypto.cpp {\tt check\_tx\_ proof()}}
    \begin{enumerate}
        \item {\em 普通地址}：\vspace{.145cm}
	    \[c' = \mathcal{H}_n(\mathfrak{m},[rK^v], [r^{resp} G + c*K^v], [r^{resp}*r G + c*k^v*r G], [T_{txprf2}], [K^v], [rG], [0])\]
	    \item {\em 子地址}：\vspace{.16cm}
	    \[c' = \mathcal{H}_n(\mathfrak{m},[rK^{v,i}], [r^{resp} K^{s,i} + c*K^{v,i}], [r^{resp}*r K^{s,i} + c*k^v*r K^{s,i}], [T_{txprf2}], [K^{v,i}], [r K^{s,i}], [K^{s,i}])\]
    \end{enumerate}{}
    \item 若 $c = c'$，则证明者知道 $k^v$，且 $k^v*r G$ 是 $K^v$ 和 $r G$ 之间的合法共享秘密（概率可忽略不计的情况除外）。
\end{enumerate}{}

\subsubsection*{使用一次性地址密钥证明「完全」所有权 (Prove `full' ownership with the one-time address key)}

虽然 InProof 表明一次性地址是用特定地址构建的（概率可忽略不计的情况除外），但这不一定意味着证明者能够{\em 花费}该 output。只有能够花费 output 的人才真正拥有它。

在 InProof 完成后，证明所有权就像用 spend key 对消息签名一样简单。\footnote{直接提供这种签名的功能在 Monero 中似乎不可用，不过正如我们将看到的，ReserveProof（第 \ref{subsec:proofs-minimum-balance-reserveproof} 节）确实包含了这种签名。}


\subsection{证明已拥有的 output 未在某交易中被花费 (UnspentProof)}
\label{subsec:proofs-owned-output-spent-unspentproof}

证明 output 已花费或未花费看似就像用 $\mathcal{J} = \{G,\mathcal{H}_p(K^o)\}$ 和 $\mathcal{K} = \{K^o,\tilde{K}\}$ 上的多基证明重新创建其 key image 一样简单。虽然这确实有效，但验证者必须知道 key image，这也会暴露未花费的 output 在{\em 未来}何时被花费。

事实证明，我们可以在不暴露 key image 的情况下证明某个 output 没有在特定交易中被花费。此外，我们可以通过将此 UnspentProof \cite{unspent-proof-issue-68} 扩展到「它作为环成员被包含在内的所有交易」来证明它{\em 目前}未被花费。\footnote{UnspentProof 尚未在 Monero 中实现。}

更具体地说，我们的 UnspentProof 说明区块链上某交易中的给定 key image 是否与其对应环中的特定一次性地址相对应。顺便提一下，正如我们将看到的，UnspentProof 与 InProof 是相辅相成的。

\subsubsection*{设置 UnspentProof}

UnspentProof 的验证者必须知道 $r K^v$，即给定已拥有的 output 的发送者-接收者共享秘密，该 output 的一次性地址为 $K^o$，交易公钥为 $r G$。他要么知道 view key $k^v$（这使他能够计算 $k^v*r G$ 并检查 $K^o - \mathcal{H}_n(k^v*rG,t)*G \stackrel{?}{=} K^s$，从而知道被测试的 output 属于证明者——回忆第 \ref{sec:one-time-addresses} 节），要么证明者提供了 $r K^v$。这就是 InProof 的用武之地，因为有了 InProof，验证者就可以确信 $r K^v$ 合法地来自证明者的 view key，并且对应一个已拥有的 output，而无需了解私有 view key。

在验证 UnspentProof 之前，验证者将知道待测试的 key image $\tilde{K}_?$，并检查其对应的环是否包含证明者所拥有 output 的一次性地址 $K^o$。然后他计算部分「花费」镜像 $\tilde{K}^s_?$。\vspace{.175cm}
\[\tilde{K}^s_? = \tilde{K}_? - \mathcal{H}_n(r K^v,t)*\mathcal{H}_p(K^o)\]

如果被测试的 key image 是从 $K^o$ 创建的，则结果点将为 $\tilde{K}^s_? = k^s*\mathcal{H}_p(K^o)$。

\subsubsection*{UnspentProof}

证明者创建两个多基证明（回忆第 \ref{sec:proofs-discrete-logarithm-multiple-bases} 节）。他的地址（拥有相关 output）为 $(K^v, K^s)$ 或 $(K^{v,i}, K^{s,i})$。\footnote{UnspentProof 对子地址和普通地址的制作方式相同。需要子地址的完整 spend key，例如 $k^{s,i} = k^s + \mathcal{H}_n(k^v,i)$（第 \ref{sec:subaddresses} 节）。}

\begin{enumerate}
    \item 一个 3 基证明，签名密钥为 $k^s$，\vspace{.175cm}
    \begin{align*}
        \mathcal{J}^{unspent}_3 &= \{[G], [K^s], [\tilde{K}^s_?]\}\\
        \mathcal{K}^{unspent}_3 &= \{[K^s], [k^s*K^s], [k^s*\tilde{K}^s_?]\}
    \end{align*}{}
    \item 一个 2 基证明，签名密钥为 $k^s*k^s$，\vspace{.175cm}
    \begin{align*}
        \mathcal{J}^{unspent}_2 &= \{[G], [\mathcal{H}_p(K^o)]\}\\
        \mathcal{K}^{unspent}_2 &= \{[k^s*K^s], [k^s*k^s*\mathcal{H}_p(K^o)]\}
    \end{align*}{}
\end{enumerate}{}

连同证明 $\sigma^{unspent}_3$ 和 $\sigma^{unspent}_2$，证明者确保传达公钥 $k^s*K^s$、$k^s*\tilde{K}^s_?$ 和 $k^s*k^s*\mathcal{H}_p(K^o)$。

\subsubsection*{验证 (Verification)}

\begin{enumerate}
    \item 确认 $\sigma^{unspent}_3$ 和 $\sigma^{unspent}_2$ 是合法的。
    \item 确保两个证明中使用了相同的公钥 $k^s*K^s$。
    \item 检查 $k^s*\tilde{K}^s_?$ 和 $k^s*k^s*\mathcal{H}_p(K^o)$ 是否相同。如果相同，则 output 已被花费；如果不同，则未被花费（概率可忽略不计的情况除外）。
\end{enumerate}{}

\subsubsection*{原理 (Why it works)}

这种看似迂回的方法防止验证者了解未花费 output 的 $k^s*H_p(K^o)$——他本可以结合 $r K^v$ 来计算其真实的 key image——同时让他确信被测试的 key image 不对应于该 output。

证明 $\sigma^{unspent}_2$ 可以用于涉及同一 output 的任意数量的 UnspentProof，尽管如果它确实被花费了，则只需要一个就足够了（即 UnspentProof 也可以用来证明一个 output 已被花费）。对给定未花费 output 被引用的所有环签名进行 UnspentProof 在计算上应该不会太昂贵。一个 output 随时间推移只可能作为诱饵被包含在大约 11 个（当前环大小）不同的环中。


\subsection{证明地址拥有最低未花费余额 ({\tt ReserveProofV2})}
\label{subsec:proofs-minimum-balance-reserveproof}

尽管在 output 未被花费时暴露其 key image 会导致隐私泄露，但这仍然是一种有些用处的方法，并在 UnspentProof 发明之前就已在 Monero 中实现 \cite{reserveproofs-pull-request-3027} \cite{unspent-proof-issue-68}。Monero 的所谓「ReserveProof」\marginnote{src/wallet/ wallet2.cpp {\tt get\_rese- rve\_proof()}}用于通过为一些未花费的 output 创建 key image 来证明某个地址拥有最低数量的资金。

更具体地说，给定一个最低余额，证明者找到足够的未花费 output 来覆盖它，用 InProof 展示所有权，为它们创建 key image 并用 2 基证明（使用不同的密钥前缀格式）证明它们是基于那些 output 合法创建的，然后用普通 Schnorr 签名证明对所用私有 spend key 的知识（如果某些 output 由不同子地址拥有，则可能有多个签名）。验证者可以检查 key image 是否未出现在区块链上，因此它们的 output 必定是未花费的。

\subsubsection*{ReserveProof}

ReserveProof 内的所有子证明签名的消息不同于其他证明（例如 OutProof、InProof 或 SpendProof）。这次是 $\mathfrak{m} = \mathcal{H}_n(\texttt{message}, \texttt{address}, \tilde{K}^o_1, ..., \tilde{K}^o_n)$，其中 {\tt address} 是证明者普通地址 $(K^v, K^s)$ 的编码形式（见 \cite{luigi-address}），key image 对应要包含在证明中的未花费 output。

\begin{enumerate}
    \item 每个 output 有一个 InProof，用于展示证明者的地址（或其某个子地址）拥有该 output。
    \item 每个 output 的 key image 用 2 基证明签名\marginnote{src/crypto/ crypto.cpp {\tt generate\_ ring\_signa- ture()}}, 挑战格式如下
    \[c = \mathcal{H}_n(\mathfrak{m}, [r G + c*K^o], [r \mathcal{H}_p(K^o) + c*\tilde{K}])\]
    \item 每个拥有至少一个 output 的地址（和子地址）都有一个普通 Schnorr 签名（第 \ref{sec:signing-messages} 节），挑战如下（普通地址和子地址相同）\marginnote{src/crypto/ crypto.cpp {\tt generate\_ signature()}}
    \[c = \mathcal{H}_n(\mathfrak{m}, K^{s,i}, [r G + c*K^{s,i}])\]
\end{enumerate}{}

将 ReserveProof 发送给其他人时\marginnote{src/wallet/ wallet2.cpp {\tt get\_rese- rve\_proof()}}[.3cm]，证明者将前缀字符串「{\tt ReserveProofV2}」与两个以 base-58 编码的列表连接起来（例如「{\tt ReserveProofV2}, list 1, list 2」）。列表 1 中的每项与特定 output 相关，包含其交易哈希（第 \ref{subsec:transaction-id} 节）、该交易中的 output 索引（第 \ref{sec:multi_out_transactions} 节）、相关共享秘密 $r K^v$、其 key image、其 InProof $\sigma^{inproof}$ 和其 key image 证明。列表 2 的项目是拥有这些 output 的地址及其 Schnorr 签名。

\subsubsection*{验证 (Verification)}

\begin{enumerate}
    \item 检查 ReserveProof 的 key image 是否未出现在区块链中。\marginnote{src/wallet/ wallet2.cpp {\tt check\_rese- rve\_proof()}}
    \item 验证每个 output 的 InProof，以及所提供的某个地址拥有每个 output。
    \item 验证 2 基 key image 签名。
    \item 使用发送者-接收者共享秘密解码 output 金额（第 \ref{sec:pedersen_monero} 节）。
    \item 检查每个地址的签名。
\end{enumerate}{}

如果一切合法，则证明者必定拥有且未花费至少 ReserveProof 的 output 中所包含的总金额（概率可忽略不计的情况除外）。\footnote{ReserveProof 虽然展示了对资金的完全所有权，但不包含证明给定子地址确实对应证明者普通地址的证明。}



\section{Monero 审计框架 (Monero audit framework)}
\label{sec:proofs-monero-audit-framework}

在美国，大多数公司每年对其财务报表进行审计 \cite{investopedia-audits}，其中包括利润表、资产负债表和现金流量表。其中前两者在很大程度上涉及公司的内部记录，而后者涉及影响公司当前持有资金的每一笔交易。加密货币是数字现金，因此对加密货币用户现金流量表的任何审计都必须与存储在区块链上的交易相关联。

被审计者的第一个任务是识别他们当前拥有的所有 output（已花费和未花费的）。这可以使用所有地址的 InProof 来完成。大型企业可能有大量的子地址，尤其是在在线市场中运营的零售商（见第 \ref{chapter:escrowed-market} 章）。为每个子地址的所有交易创建 InProof 可能会导致证明者和验证者产生巨大的计算和存储需求。

相反，我们可以仅为证明者的普通地址（对所有交易）制作 InProof。审计者使用那些发送者-接收者共享秘密来检查是否有任何 output 由证明者的主地址或其相关子地址拥有。回忆第 \ref{sec:subaddresses} 节，用户的 view key 足以识别地址的子地址拥有的所有 output。

为了确保证明者不会通过隐藏某些子地址的普通地址来欺骗审计者，他还必须证明所有子地址都对应于他已知的某个普通地址。


\subsection{证明地址和子地址的对应关系 (SubaddressProof)}
\label{subsec:proofs-address-subaddress-correspond-subaddressproof}

SubaddressProof 展示普通地址的 view key 可用于识别给定子地址拥有的 output。\footnote{SubaddressProof 尚未在 Monero 中实现。}

\subsubsection*{SubaddressProof}
%subaddress proof: base keys [G, K^{s,i}] public keys [K^v, K^{v,i}] signing key [k^v]
SubaddressProof 的制作方式与 OutProof 和 InProof 大致相同。此处基密钥为 $\mathcal{J} = \{G, K^{s,i}\}$，公钥为 $\mathcal{K} = \{K^v, K^{v,i}\}$，签名密钥为 $k^v$。同样，我们仅展示验证步骤以说明我们的意思。

\subsubsection*{验证 (Verification)}

验证者知道证明者的地址 $(K^v, K^s)$、子地址 $(K^{v,i}, K^{s,i})$，并拥有 SubaddressProof $\sigma^{subproof} = (c,r)$。

\begin{enumerate}
    \item 计算挑战\vspace{.175cm}
	\[c' = \mathcal{H}_n(\mathfrak{m},[K^{v,i}], [r G + c*K^v], [r K^{s,i} + c*K^{v,i}], [T_{txprf2}], [K^v], [K^{s,i}], [0])\]
    \item 若 $c = c'$，则证明者知道 $K^v$ 的 $k^v$，且 $K^{s,i}$ 与该 view key 结合可以生成 $K^{v,i}$（概率可忽略不计的情况除外）。
\end{enumerate}{}


\subsection{审计框架 (The audit framework)}
\label{subsec:audit-framework}

现在我们已经准备好尽可能多地了解一个人的交易历史。\footnote{此审计框架在 Monero 中尚未完全可用。SubaddressProof 和 UnspentProof 尚未实现，InProof 尚未为我们所解释的子地址相关优化做好准备，并且没有真正的结构来方便地获取或组织证明者和验证者所需的所有信息。}

\begin{enumerate}
    \item 证明者收集他所有账户的列表，每个账户由一个普通地址和各种子地址组成。他为所有子地址制作 SubaddressProof。与 ReserveProof 类似，他还为每个地址和子地址的 spend key 制作签名，展示他对这些地址所拥有的所有 output 的花费权。
    \item 证明者为他的每个普通地址在区块链上的所有交易（即所有交易公钥）生成 InProof。这向审计者揭示了证明者的地址拥有的所有 output，因为他们可以用发送者-接收者共享秘密检查所有一次性地址。他们可以确信子地址拥有的 output 也会被识别，因为有 SubaddressProof。\footnote{此步骤也可以通过提供私有 view key 来完成，但这有明显的隐私影响。}
    \item 证明者为他的每个已拥有的 output，在它们作为环成员出现的所有交易 input 上生成 UnspentProof。现在审计者将知道证明者的余额，并可以进一步调查已花费的 output。\footnote{或者，他可以为所有已拥有的 output 制作 ReserveProof。同样，暴露未花费 output 的 key image 有明显的隐私影响。}
    \item {\em 可选}：证明者为他花费 output 的每笔交易生成 OutProof，向审计者展示接收者和金额。此步骤仅适用于证明者保存了交易私钥的交易。
\end{enumerate}{}

重要的是，证明者无法直接展示资金来源。他唯一的办法是向给他汇款的人请求一组证明。

\begin{enumerate}
    \item 对于向证明者汇款的交易，其作者制作 SpendProof 证明确实发送了它。
    \item 证明者的资助者还用一个标识性公钥（例如其普通地址的 spend key）制作签名。SpendProof 和这个签名都应签名包含该标识性公钥的消息，以确保 SpendProof 没有被窃取，也不是由其他人制作的。
\end{enumerate}{}
